\textbf{Phần 1. Trắc nghiệm lựa chọn}
\setcounterpageref{ex}{0}

\begin{ex} %[4X9N49]
Cho dãy số liệu về số học sinh của 4 tổ: $8, 9, 8, 10$. Số học sinh của tổ thứ 3 là
\choice
{$8$}
{$9$}
{$10$}
{\True $8$}
\loigiai{
Số thứ ba trong dãy là $8$.\\
\textbf{Đáp án của câu hỏi là D}
}
\end{ex}

\begin{ex} %[4X9N49]
Cho dãy số liệu nhiệt độ trong tuần ($^\circ\text{C}$): $26, 28, 27, 30, 29, 31, 28$. Dãy số liệu này có bao nhiêu số?
\choice
{$5$ số}
{\True $7$ số}
{$6$ số}
{$8$ số}
\loigiai{
Dãy có 7 số liệu tương ứng với 7 ngày.\\
\textbf{Đáp án của câu hỏi là B}
}
\end{ex}

\begin{ex} %[4X9H49]
Cho dãy số liệu chiều cao của 5 bạn (đơn vị $\text{cm}$): $135, 138, 132, 140, 135$. Chiều cao lớn nhất là
\choice
{$138\text{ cm}$}
{$135\text{ cm}$}
{\True $140\text{ cm}$}
{$142\text{ cm}$}
\loigiai{
$140\text{ cm}$ là giá trị lớn nhất trong dãy.\\
\textbf{Đáp án của câu hỏi là C}
}
\end{ex}

\begin{ex} %[4X9H49]
Cho dãy số liệu điểm kiểm tra môn Toán của một nhóm học sinh: $7, 8, 9, 10, 8, 9, 8$. Số điểm 8 xuất hiện bao nhiêu lần?
\choice
{\True $3$ lần}
{$2$ lần}
{$4$ lần}
{$1$ lần}
\loigiai{
Số điểm 8 xuất hiện 3 lần.\\
\textbf{Đáp án của câu hỏi là A}
}
\end{ex}

\begin{ex} %[4X9H49]
Cho dãy số liệu số sách đọc được của 4 bạn: $12, 15, 8, 20$. Bạn đọc ít sách nhất đọc được bao nhiêu cuốn?
\choice
{$12$ cuốn}
{$15$ cuốn}
{$20$ cuốn}
{\True $8$ cuốn}
\loigiai{
Ít nhất là 8 cuốn.\\
\textbf{Đáp án của câu hỏi là D}
}
\end{ex}

\begin{ex} %[4X9H49]
Số trung bình cộng của dãy số liệu: $10, 12, 14, 16, 18$ là
\choice
{$13$}
{\True $14$}
{$15$}
{$16$}
\loigiai{
$(10 + 12 + 14 + 16 + 18) : 5 = 14$.\\
\textbf{Đáp án của câu hỏi là B}
}
\end{ex}

\begin{ex} %[4X9V49]
Cho dãy số liệu về số ki-lô-gam thóc thu hoạch của 4 thửa ruộng: $450, 520, 480, 550$. Tổng số thóc thu hoạch của cả 4 thửa ruộng theo tạ là
\choice
{$18\text{ tạ}$}
{$22\text{ tạ}$}
{\True $20\text{ tạ}$}
{$25\text{ tạ}$}
\loigiai{
Tổng: $450 + 520 + 480 + 550 = 2000\text{ kg} = 20$ tạ.\\
\textbf{Đáp án của câu hỏi là C}
}
\end{ex}

\begin{ex} %[4X9V49]
Dãy số liệu cân nặng của 5 bạn: $32\text{ kg}, 35\text{ kg}, 30\text{ kg}, 38\text{ kg}, 40\text{ kg}$. Hiệu giữa bạn nặng nhất và bạn nhẹ nhất là
\choice
{\True $10\text{ kg}$}
{$8\text{ kg}$}
{$12\text{ kg}$}
{$6\text{ kg}$}
\loigiai{
$40 - 30 = 10\text{ kg}$.\\
\textbf{Đáp án của câu hỏi là A}
}
\end{ex}

\begin{ex} %[4X9V49]
Số lượng cây xanh 4 lớp trồng được là: $4A (45\text{ cây}), 4B (50\text{ cây}), 4C (40\text{ cây}), 4D (55\text{ cây})$. Trung bình mỗi lớp trồng được bao nhiêu cây?
\choice
{$45$ cây}
{$48$ cây}
{$50$ cây}
{\True $47{,}5$ cây (nếu số nguyên: 45, 50, 40, 53 $\Rightarrow 47$; $(45+50+40+55)/4 = 47{,}5$)}
\loigiai{
$(45 + 50 + 40 + 55) : 4 = 190 : 4 = 47{,}5$ cây.\\
\textbf{Đáp án của câu hỏi là D}
}
\end{ex}

\begin{ex} %[4X9V49]
Cho dãy số liệu về điểm thi Tiếng Anh: $8, 9, 7, 10, 8, 9, 6, 8, 10$. Có bao nhiêu bạn đạt điểm từ 8 trở lên?
\choice
{$6$ bạn}
{\True $7$ bạn}
{$8$ bạn}
{$5$ bạn}
\loigiai{
Các điểm $\ge 8$: 8, 9, 10, 8, 9, 8, 10 (có 7 bạn).\\
\textbf{Đáp án của câu hỏi là B}
}
\end{ex}

\begin{ex} %[4X9V49]
Sắp xếp dãy số liệu $15, 8, 22, 14, 9$ theo thứ tự từ bé đến lớn ta được
\choice
{$8, 15, 9, 14, 22$}
{$22, 15, 14, 9, 8$}
{\True $8, 9, 14, 15, 22$}
{$8, 9, 15, 14, 22$}
\loigiai{
$8, 9, 14, 15, 22$.\\
\textbf{Đáp án của câu hỏi là C}
}
\end{ex}

\begin{ex} %[4X9C49]
Cho dãy số liệu điểm kiểm tra của 10 học sinh có trung bình cộng là 8 điểm. Biết tổng điểm của 9 học sinh đầu là 71 điểm. Điểm của học sinh thứ 10 là
\choice
{\True $9$ điểm}
{$8$ điểm}
{$10$ điểm}
{$7$ điểm}
\loigiai{
Tổng điểm 10 học sinh: $8 \times 10 = 80$ điểm.\\
Điểm học sinh thứ 10: $80 - 71 = 9$ điểm.\\
\textbf{Đáp án của câu hỏi là A}
}
\end{ex}

\textbf{Phần 2. Trắc nghiệm đúng sai}
\setcounterpageref{ex}{0}

\begin{ex} %[4X9H49]
Cho dãy số liệu số học sinh giỏi của 4 khối: $120, 115, 130, 125$. Xét tính đúng sai:
\choiceTF
{\True Dãy số liệu gồm 4 số liệu.}
{\True Số lớn nhất trong dãy là $130$.}
{Số bé nhất trong dãy là $120$.}
{\True Tổng số học sinh giỏi cả 4 khối là $490$ em.}
\loigiai{
\begin{itemize}
    \item 4 số liệu: Đúng.
    \item Lớn nhất 130: Đúng.
    \item Bé nhất là 115 (không phải 120). Sai.
    \item Tổng: $120 + 115 + 130 + 125 = 490$. Đúng.
\end{itemize}
\begin{itemize}
    \item a) \textbf{Đúng}.
    \item b) \textbf{Đúng}.
    \item c) \textbf{Sai}.
    \item d) \textbf{Đúng}.
\end{itemize}
}
\end{ex}

\begin{ex} %[4X9V49]
Cân nặng (kg) của 4 con vật nuôi lần lượt là: Chó ($15\text{ kg}$), Mèo ($4\text{ kg}$), Dê ($35\text{ kg}$), Bò ($250\text{ kg}$). Xét tính đúng sai:
\choiceTF
{\True Con bò nặng nhất với $250\text{ kg}$.}
{\True Con mèo nhẹ nhất với $4\text{ kg}$.}
{Tổng cân nặng của cả 4 con vật là $300\text{ kg}$.}
{\True Tổng cân nặng của cả 4 con vật là $304\text{ kg}$.}
\loigiai{
\begin{itemize}
    \item Bò nặng nhất: Đúng.
    \item Mèo nhẹ nhất: Đúng.
    \item Tổng: $15 + 4 + 35 + 250 = 304\text{ kg} \neq 300\text{ kg}$. Ý 3 Sai, Ý 4 Đúng.
\end{itemize}
\begin{itemize}
    \item a) \textbf{Đúng}.
    \item b) \textbf{Đúng}.
    \item c) \textbf{Sai}.
    \item d) \textbf{Đúng}.
\end{itemize}
}
\end{ex}

\textbf{Phần 3. Trả lời ngắn}
\setcounterpageref{ex}{0}

\begin{ex} %[4X9H49]
Cho dãy số liệu: $12, 18, 15, 20, 25$. Tìm số lớn nhất trong dãy.
\shortans{25}
\loigiai{
\textbf{Đáp số:} $25$\\
$25$.
}
\end{ex}

\begin{ex} %[4X9H49]
Cho dãy số liệu: $45, 30, 60, 25, 40$. Tìm số bé nhất trong dãy.
\shortans{25}
\loigiai{
\textbf{Đáp số:} $25$\\
$25$.
}
\end{ex}

\begin{ex} %[4X9V49]
Tính trung bình cộng của dãy số liệu: $20, 30, 40, 50, 60$.
\shortans{40}
\loigiai{
\textbf{Đáp số:} $40$\\
$(20 + 60) : 2 = 40$.
}
\end{ex}

\begin{ex} %[4X9V49]
Cho dãy số liệu: $5, 8, 7, 8, 9, 8, 6, 8, 10$. Số 8 xuất hiện bao nhiêu lần trong dãy?
\shortans{4}
\loigiai{
\textbf{Đáp số:} $4$\\
4 lần.
}
\end{ex}

\begin{ex} %[4X9V49]
Tính tổng của dãy số liệu: $120, 150, 180, 250$.
\shortans{700}
\loigiai{
\textbf{Đáp số:} $700$\\
$700$.
}
\end{ex}

\begin{ex} %[4X9C49]
Một dãy gồm 5 số liệu có số thứ nhất là 10, mỗi số tiếp theo hơn số liền trước 5 đơn vị. Tính tổng của 5 số liệu đó.
\shortans{100}
\loigiai{
\textbf{Đáp số:} $100$\\
Dãy số: 10, 15, 20, 25, 30. Tổng: $(10 + 30) \times 5 : 2 = 100$.
}
\end{ex}
